\section{Introduction} \label{section:intro} 

Let $p$ be a prime, and let $K_\infty$ be a $\Z_p$-extension of a number field $K$. Then for each~${m \in \N}$ there exists a unique subextension $K_m/K$ of degree $p^m$. Let $h_m$ be the class number of $K_m$, and let ${v_p(h_m) \ge 0}$ be the exponent of the largest power of~$p$ which divides $h_m$, respectively. Then Iwasawa (see~\cite{iwasawa}) proved that we have an asymptotic formula 
\begin{align} \label{eq:iw} 
	v_p(h_m) = \mu(K_\infty/K) \cdot p^m + \lambda(K_\infty/K) \cdot m + \nu(K_\infty/K) 
\end{align} 
for each sufficiently large $m$, where $\mu(K_\infty/K)\in \Z_{\ge 0}$, $\lambda(K_\infty/K)\in \Z_{\ge 0}$ and $\nu(K_\infty/K)\in \Z$ denote the Iwasawa invariants of the $\Z_p$-extension $K_\infty/K$. More generally, let $K_\infty$ be a $\Z_p^l$-extension of a number field $K$, with intermediate fields $K_m$ (i.e. $K_m$ is the unique intermediate field such that ${\Gal(K_m/K) \cong (\Z/p^m\Z)^l}$, ${m \in \N}$). Again, we denote by $h_m$ the class number of $K_m$. Then Greenberg conjectured (see ~\cite[Section~7]{cuoco-monsky}) that there exists a polynomial ${P(X,Y) \in \Q[X,Y]}$ of total degree at most~$l$ and of degree at most 1 in $Y$ such that 
\[ v_p(h_m) = P(p^m,m) \] 
for each sufficiently large $m$. In this paper, we will refer to this conjecture as \emph{Greenberg's conjecture}. Cuoco and Monsky have generalised Iwasawa's formula to $\Z_p^l$-extensions, but their formula contains an error term $O(p^{n(d-1)})$ and so they were not able to prove Greenberg's conjecture (in fact, it seems that this conjecture might not even be true in general, see~\cite{cuoco-monsky}). 

In the Iwasawa theory of graphs one starts with a finite connected graph $X$ instead of a number field $K$ and one considers a sequence of Galois covers
\[X\leftarrow X_1\leftarrow X_2\leftarrow X_3\leftarrow X_4\leftarrow \cdots \] 
such that $\Gal(X_m/X)\cong (\Z/p^m\Z)^l$. Then each $X_m$ is again a finite graph and all the $X_m$ are connected. Instead of the $p$-part of the class number of the field $K_m$ one considers the $q$-part of the number of spanning trees of $X_m$ (where $q$ is a prime not necessarily distinct from $p$). Valli\`eres and McGown-Valli\`eres proved in a sequence of papers (see~\cite{vallieres1},~\cite{vallieres2},~\cite{vallieres3}) that (under mild technical hypothesis) the number of spanning trees in a $\Z_p$-cover (i.e. ${l = 1}$) grows asymptotically as in Iwasawa's formula. In a subsequent work Lei and Valli\`eres considered the case ${q \neq p}$ (see~\cite{vallieres-lei}). The proofs of the above results are purely analytic and make use of Ihara $L$-functions. On the other hand, Gonet proved an Iwasawa-like formula in the case ${p = q}$ via a more algebraic approach, which is closer to Iwasawa's module theoretical approach (see~\cite{gonet}). 

In this paper we study analogous growth patterns in the theory of voltage covers of graphs for $l\ge2$ and ${q =p}$. Like Gonet's one our approach is purely algebraic and uses the machinery developed by Cuoco and Monsky, together with the fact that the number of spanning trees is given by the size of the Jacobian of a finite connected graph (see Section~\ref{section:notation} for precise definitions). We are able to prove Greenberg's conjecture under the assumption that all the involved graphs are connected (see Theorem~\ref{thm:greenberg-conj}). This result had also been obtained before under some mild additional assumptions by Valli\`eres and DuBose, generalising the approach of McGown-Valli\`eres, using Ihara $L$-functions ~\cite[Theorem~A]{dubose-vallieres}. 

A central aspect of classical Iwasawa theory is the interplay between algebraic objects such as certain Galois modules on the one hand and analytic objects such as $L$-functions (e.g. the Hasse-Weil $L$-function $L(E,s)$ for an elliptic curve $E$ or the Dedekind zeta function $\zeta_K(s)$ for a number field $K$) on the other hand. In many cases one can define a $p$-adic function interpolating special values of these complex analytic functions. Given a finite connected graph $X$ and a finite Galois cover $Y/X$ we consider the complex analytic Ihara $L$-function $L_X(\chi,s)$ for every character ${\chi\in \widehat{\Gal(Y/X)}}$. This $L$-function has a factor $(1-s^2)^{\chi(X)}$, where $\chi(X)$ is the Euler characteristic of~$X$. As this factor is independent of the cover $Y/X$ and just produces trivial zeros at $s=1$, we will ignore it and view ${Q_X(\chi,s)=L_X(\chi,s)/(1-s^2)^{\chi(X)}}$ as the algebraic part of the Ihara $L$-function. 

Let now $X$ be a finite connected graph (we allow multiple edges and loops), and suppose that $X$ has $n$ vertices. Let $X_m$ be the $m$-th intermediate graph in a $\Z_p^l$-cover $X_\infty$ of $X$, i.e. ${\Gal(X_m/X) \cong (\Z/p^m\Z)^l}$. Then there exists an element 
\[ \Delta_\infty\in \textup{Mat}_{n,n}(\Z_p[[\Gal(X_\infty/X)]])\] 
such that $\chi(\det(\Delta_\infty))$ interpolates the algebraic part of $L_X(\chi,s)$ at $s=1$ in the sense that \[\chi(\det(\Delta_\infty))^{-1}=Q_X(\chi,1).\] In this setting we obtain the following 
\begin{theo*}[Iwasawa main conjecture, Theorem~\ref{thm:main_conjecture}]
	Let $J(X_\infty)$ be the Jacobian of $X_\infty$ and let ${R = \Z_p[T_1, \ldots, T_l]}$ and ${\Lambda=\Z_p\llbracket T_1, \ldots, T_l\rrbracket }$. Then ${J(X_\infty) \otimes_R \Lambda}$ is a finitely generated torsion $\Lambda$-module, and 
	\[\textup{Char}_\Lambda(J(X_\infty)\otimes_R \Lambda)=(\det(\Delta_\infty)).\]
\end{theo*}
Here $\textup{Char}_\Lambda(M)$ denotes the characteristic ideal of a finitely generated and torsion $\Lambda$-module $M$ (for details we refer to Section~\ref{section:notation_iw-modules}). This should be understood as an analogous statement to the classical Iwasawa main conjecture along $\Q_p(\mu_{p^\infty})$ that relates the characteristic ideals of the projective limits of $p$-class groups to power series interpolating special values of Dirichlet $L$-functions. 

The last central topic which we want to address in the present paper is the behaviour of so-called \emph{(generalised) Iwasawa invariants} if we vary the $\Z_p^l$-cover of a given finite connected graph $X$. Let $P(X,Y)$ be the polynomial appearing in Greenberg's conjecture. Then
\[P(p^m,m)=m_0p^{ml}+l_0mp^{m(l-1)}+O(p^{m(l-1)}).\]
Following Cuoco and Monsky (see~\cite{cuoco-monsky}) we call $m_0$ and $l_0$ the (generalised) Iwasawa invariants of $J(X_\infty)\otimes \Lambda$, and we write $m_0(X_\infty)$ and $l_0(X_\infty)$ for these parameters. In Section~\ref{sec:Fukuda} we define a topology on the set $\mathcal{E}^l(X)$ of voltage $\Z_p^l$-covers of $X$ and we prove the following results. 
\begin{theo*}[Theorems~\ref{thm:local_max1} and~\ref{thm:local_max2}]
	We have the following two results. 
	\begin{compactenum}
		\item Assume that $l=1$ and let ${X_\infty \in \mathcal{E}^1(X)}$. Then there exists a (sufficiently small) neighbourhood $U$ of $X_\infty$ such that the following statements hold: \begin{compactenum}[(a)] 
			\item For each ${\tilde{X}_\infty \in U}$, we have 
			\[ m_0(\tilde{X}_\infty) \le m_0(X_\infty). \] 
			\item ${l_0(\tilde{X}_\infty) = l_0(X_\infty)}$ for each ${\tilde{X}_\infty \in U}$ for which ${m_0(\tilde{X}_\infty) = m_0(X_\infty)}$. 
							\end{compactenum} 
		\item Fix an element ${X_\infty \in \mathcal{E}^l(X)}$, ${l \ge 2}$. Then there exist an integer ${k \in \N}$ and a neighbourhood $U$ of $X_\infty$ such that the following two statements hold for each ${\tilde{X}_\infty}$ in $U$. \begin{compactenum}[(a)] 
			\item $m_0(\tilde{X}_\infty) \le m_0(X_\infty)$, and 
			\item $l_0(\tilde{X}_\infty) \le k$ holds if ${m_0(\tilde{X}_\infty) = m_0(X_\infty)}$. 
		\end{compactenum} 
	\end{compactenum}
\end{theo*}

Let us briefly describe the structure of the article. Section~\ref{section:notation} is preliminary in nature -- here we introduce the basic notation and the setup. In Section~\ref{section:J(X_m)} we relate the Jacobians of the Galois covers $X_m/X$, ${m \in \N}$, to certain quotients of an Iwasawa module. This enables us to describe the growth of these Jacobians by using results of Cuoco and Monsky, and we already obtain a weak result in the direction of Greenberg's conjecture (see Lemma~\ref{lemma:NundL}). In Section~\ref{section:linearalgebra}, we prove Greenberg's conjecture by studying the voltage $p$-Laplacian of the Galois cover and using a suitable matrix representation in order to describe the growth of the Jacobians in terms of a certain power series. In Section~\ref{section:main_conjecture} we prove the Iwasawa main conjecture and in Section~\ref{sec:Fukuda} we prove Theorems~\ref{thm:local_max1} and~\ref{thm:local_max2} which have been stated above. The proofs are completely algebraic, and they are based on work of Fukuda (see~\cite{fukuda}) and the first named author (see~\cite{local_beh} and~\cite{local_max}). In the final two sections, we compute many numerical examples by applying the results from Sections~\ref{section:linearalgebra} and ~\ref{section:main_conjecture}. In particular, we construct examples of Galois covers $X_\infty/X$ of graphs with non-trivial $m_0$- and $l_0$-invariant. To construct these examples we need to be able to compare the Jacobians in a $\Z_p^l$-cover $X_\infty/X$ with the Jacobian of a $\Z_p$-cover $Y/X$ contained in $X_\infty$. In the setting of Iwasawa theory of class groups or elliptic curves results of this form are often referred to as control theorems. We will prove the control theorem needed in our context in Section~\ref{sec:control}. 

In the classical Iwasawa theory of $\Z_p$-extensions (i.e. ${l = 1}$) of number fields, Iwasawa himself was the first who constructed examples with a non-trivial $m_0$-invariant (see~\cite{iwasawa_mu}). On the other hand, to the authors' knowledge no example of a $\Z_p^l$-extension (for $l\ge 2$) of a number field $K$ is known where one can show that the $l_0$-invariant is greater than zero (however, there exist such examples in the setting of Selmer groups of elliptic curves, see~\cite{kleine-matar}). Therefore the construction of an example with non-trivial $l_0$-invariant in the final section might be of some interest. 

\section{Notation and definitions} \label{section:notation} 
\subsection{Iwasawa modules} \label{section:notation_iw-modules} 
Let $\Gamma$ be a group which is topologically isomorphic to $\Z_p^l$. Then the completed group ring $\Z_p\llbracket \Gamma\rrbracket $ is defined as the projective limit 
\[ \Z_p\llbracket \Gamma\rrbracket = \varprojlim_{U} \Z_p[\Gamma/U],\] 
where ${U \subseteq \Gamma}$ runs over the open normal subgroups. The ring $\Z_p\llbracket \Gamma\rrbracket $ can be identified (non-canonically) with the ring $\Z_p\llbracket T_1, \ldots, T_l\rrbracket $ of formal power series in $l$ variables, as follows: choose a set of topological generators ${\gamma_1, \ldots, \gamma_l}$ of $\Gamma$, and consider the bijective homomorphism between the group rings which is induced by mapping $\gamma_i$ to ${T_i + 1}$, respectively. 

In the following we always denote by $\Lambda_l = \Z_p\llbracket T_1, \ldots, T_l \rrbracket $ the Iwasawa algebra in $l$ variables. If $l$ is clear from the context, then we abbreviate $\Lambda_l$ to $\Lambda$. 

Now we describe the structure theory of \emph{Iwasawa modules}. In this article, any finitely generated $\Lambda$-module $A$ will be called an Iwasawa module. An Iwasawa module is called \emph{pseudo-null} if it is annihilated by two relatively prime elements of the unique factorisation domain $\Lambda$. A \emph{pseudo-isomorphism} of finitely generated $\Lambda$-modules $A$ and $B$ is a $\Lambda$-module homomorphism ${\varphi \colon A \longrightarrow B}$ such that the kernel and the cokernel of $\varphi$ are pseudo-null. 

Let now $A$ be any Iwasawa module. Then by the general structure theory (see ~\cite[Section~5.1]{nsw}) there exist an \emph{elementary} $\Lambda$-module 
\[ E_A = \Lambda^r \oplus \bigoplus_{i = 1}^s \Lambda/(p^{e_i}) \oplus \bigoplus_{j = 1}^t \Lambda/(g_j) \] 
and a pseudo-isomorphism ${\varphi \colon A \longrightarrow E_A}$. Here $g_1, \ldots, g_t$ are elements of $\Lambda$ which are coprime with $p$. If ${r = 0}$, then the \emph{characteristic power series} $F_A$ attached to $A$ is the element 
\[ F_A = p^{\sum_{i = 1}^s e_i} \cdot \prod_{j = 1}^t g_j, \] 
and the \emph{characteristic ideal} $\textup{Char}(E_A)$ of $E_A$  (and $A$) is the principal ideal $(F_A)$. If ${r \ge 1}$, i.e. if $A$ is not $\Lambda$-torsion, then we define ${F_A = 0}$ and ${\textup{Char}(A) = (0)}$. The characteristic power series is determined up to units of $\Lambda$ by the Iwasawa module $A$, and therefore the characteristic ideal is well-defined. 

Now suppose that $A$ is $\Lambda$-torsion. Then we define the \emph{(generalised) Iwasawa invariants} of $A$ as follows. To this purpose, we fix an isomorphism between $\Z_p\llbracket \Gamma\rrbracket $ and~$\Lambda$, as above, and we consider the characteristic power series $F_A$ attached to $A$. Let 
\[ m_0(A) = \sum_{i = 1}^s e_i, \] 
i.e. $p^{m_0(A)}$ is the exact power of $p$ which divides $F_A$. Moreover, we write 
\[ F_A = p^{m_0(A)} \cdot G_A, \] 
i.e. $G_A$ is coprime with $p$, and we consider the coset $\overline{G_A}$ of $G_A$ in the quotient algebra $\overline{\Lambda} = \Lambda/(p)$. Then we define 
\[ l_0(A) = \sum_{\mathcal{P}} v_{\mathcal{P}}(\overline{G_A}), \] 
where the sum runs over all the prime ideals of $\overline{\Lambda}$ of the form ${\mathcal{P} = (\overline{\sigma - 1})}$ for some element ${\sigma \in \Gamma \setminus \Gamma^p}$, and where $v_{\mathcal{P}}$ denotes the $\mathcal{P}$-adic valuation. Note that this sum is finite since $\overline{\Lambda}$ is again a unique factorisation domain. 

In the case of ${l = 1}$, the structure theory of Iwasawa modules is better understood. In this case, an Iwasawa module is pseudo-null if and only if it is \emph{finite}. Moreover, it follows from the Weierstra{\ss} Preparation Theorem for ${\Lambda = \Z_p\llbracket T\rrbracket }$ (see ~\cite[Theorem~7.3]{wash}) that the characteristic power series $F_A$ of a torsion $\Lambda$-module $A$ is associated to a \emph{polynomial} of the form 
\[ p^{m_0(A)} \cdot G_A, \] 
where $G_A$ is \emph{distinguished}. This means that $G_A$ is monic, and that each coefficient of~$G_A$, apart from the leading one, is divisible by $p$. Then we define the \emph{classical Iwasawa invariants} of $A$ as 
\[ \mu(A) = m_0(A), \quad \lambda(A) = \deg(G_A)=l_0(G_A). \] 
If $A = \varprojlim A_n$ is the classical Iwasawa module of ideal class groups in the intermediate layers of a $\Z_p$-extension $K_\infty/K$, then the Iwasawa invariants of $A$ are the constants from the asymptotic growth formula~\eqref{eq:iw} mentioned in the introduction (see~\cite{iwasawa} or ~\cite[Chapter~13]{wash} for more details). 

We conclude the current section by introducing some more notation. For any ${i \in \{1, \ldots, l\}}$ and every ${m \in \N}$, we let 
\[ \omega_m(T_i) = (T_i + 1)^{p^m} - 1. \] 
Moreover, for every $0 \le j \le m$, we introduce 
\[ \nu_{m,j}(T_i) =\omega_m(T_i)/\omega_j(T_i). \] 
Recall that $\Lambda \cong \Z_p[[\Gamma]]$, where $\Gamma$ is topologically isomorphic to $\Z_p^l$. In fact, the polynomials $\omega_m(T_i)$ correspond to elements in ${R = \Z_p[\Gamma] \subseteq \Lambda}$. We denote by $I_m$ the ideal of ${\Lambda}$ generated by all the $\omega_m(T_i)$, and we denote the quotient ${\Lambda/I_m}$ by $R_m$, ${m \in \N}$. Note that ${R_m \cong \Z_p[\Gamma/\Gamma^{p^m}]}$. 

\subsection{Graphs} 
We briefly define the notions from graph theory which will be needed; for a more detailed introduction see for example the paper~\cite{dubose-vallieres} of DuBose and Valli\`eres. 

Let $Y$ be a graph, and let $\mathbf{E}_Y$ and $E_Y$ be the sets of \emph{directed} and \emph{undirected edges} of~$Y$. If ${e \in \mathbf{E}_Y}$ has origin~$v$ and target $w$ then we denote by $e^{-1}$ the directed edge with origin $w$ and target $v$. 
In this article we denote the set of directed edges between two vertices $v$ and $w$ of a graph $Y$ by ${E(v,w) \subseteq \mathbf{E}_Y}$, i.e. each ${e \in E(v,w)}$ has origin~$v$ and target $w$. Sometimes we say that ${e \in E(v,w)}$ is an edge \emph{from $v$ to $w$}. Note that~$E(v,w)$ may contain more than one element (i.e. we allow \emph{multigraphs}). For any graph~$Y$ we make the following definitions. 
\begin{defi}
	Let $Y$ be a (not necessarily finite) graph. Then we denote by $V(Y)$ the set of vertices of $Y$. We let 
	\[ \deg(v) = \sum_{w \in V(Y), w\neq v} |E(v,w)|+2\vert \{\textup{undirected loops from $v$ to $v$}\}\vert\]
	be the \emph{degree} of a vertex $v$ and we define $\textup{mult}(v,w) = |E(v,w)|$ to be the number of edges with origin $v$ and target $w$ (i.e. $\textup{mult}(v,w) = 0$ if $v$ and $w$ are not adjacent). 
	
	In this article we always assume that $\deg(v)$ is finite for each ${v \in V(Y)}$ (such graphs are called locally finite). 
	
	We define
	\begin{align*}
		\textup{Div}(Y) &=\left\{\sum_{v \in V(Y)} a_v v \; \Big| \; a_v\in \Z_p, a_v=0 \textup{ for all but finitely many }v\right\}\\
		\textup{Div}^0(Y) &=\left\{d = \sum_{v \in V(Y)} a_v v \in \textup{Div}(Y) \; \Big| \; \sum_{v \in V(Y)} a_v=0\right\} \\ 
		\textup{Pr}(Y) &=\mathcal{L}(\textup{Div}(Y)), 
	\end{align*}
	where $\mathcal{L}$ is the Laplacian operator which is defined as follows: For a vertex ${v\in V(Y)}$ we define ${\mathcal{D}(v)=\deg(v)v}$ and ${\mathcal{A}(v)=\sum_{w\sim v} \textup{mult}(v,w) w}$, where ${w \sim v}$ means that the vertices $w$ and $v$ are adjacent. We now define \[\mathcal{L}(v)=\mathcal{D}(v)-\mathcal{A}(v).\]
	We further define the two quotient groups 
	\begin{align*}
		\textup{Pic}_p(Y) &=\textup{Div}(Y)/\textup{Pr}(Y), \\
		J_p(Y) &=\textup{Div}^0(Y)/\textup{Pr}(Y).
	\end{align*}
\end{defi}
The subscript $p$ is used to emphasize that we are considering $\Z_p$-modules here. Let $\textup{Jac}(Y)$ be the \emph{Jacobian} of $Y$. Then $J_p(Y)=\textup{Jac}(Y)\otimes \Z_p$. Note that $J_p(Y)$ is the $p$-Sylow subgroup of $\textup{Jac}(Y)$ if $\textup{Jac}(Y)$ is finite. As the prime $p$ is fixed once and for all we will write $J(Y)$ instead of $J_p(Y)$ from now on. If $Y$ is finite and connected, then it is well-known that $J(Y)$ is a finite abelian group (see ~\cite[Proposition 2.37]{corry-Perkinson}). Furthermore, the group $\textup{Pr}(Y)$ is generated by the elements
\[p_v= \mathcal{L}(v) = \deg(v)v-\sum_{w\sim v, w\neq v}\textup{mult}(v,w)w-2\vert \{\textup{undirected loops from $v$ to $v$}\}\vert v, \] 
$v \in V(Y)$. 

Let now $X$ be a finite graph with $n$ vertices. We specialise to the situation of voltage covers $X_m$ (see also~\cite{diss-gonet}). 
Let $\gamma\colon E_X\longrightarrow\mathbf{E}_X$ be a section of the natural map $\mathbf{E}_X\longrightarrow E_X$. Let ${S=\gamma(E_X)}$. The set ${S \subseteq \mathbf{E}_X}$ corresponds to a choice of orientation. 

Let $\Gamma$ be a multiplicative group which is isomorphic to $\Z_p^l$. We fix a \emph{voltage assignment}, i.e. a map  
\[\alpha \colon S\longrightarrow \Gamma.\]
Using the natural projection maps ${\Gamma\longrightarrow \Gamma/\Gamma^{p^m}}$, ${m \in \N}$, $\alpha$ induces well-defined assignments
\[\alpha_m\colon S\longrightarrow \Gamma/\Gamma^{p^m}.\]
Since the target groups of $\alpha$ in our applications usually will be Galois groups, we write the target groups of $\alpha$ and of the $\alpha_m$ multiplicatively. Note that for each element $e\in \mathbf{E}_X$ either $e\in S$ or $e^{-1}\in S$. If $e\notin S$, we define $\alpha(e)=\alpha(e^{-1})^{-1}$ and analogously for $\alpha_m$. Thus, we can interpret $\alpha$ and $\alpha_m$ as maps defined on $\mathbf{E}_X$. 
\begin{defi} \label{def:voltagelevel} 
	We define the derived graph ${X_m := X(\Gamma/\Gamma^{p^m},S,\alpha_m)}$ on the vertices $(v,g\!\!\pmod{\Gamma^{p^m}})$, where $g\in \Gamma$ and $g \!\!\pmod{\Gamma^{p^m}}$ denotes the equivalence class of $g$ in~$\Gamma/\Gamma^{p^m}$. The set $\mathbf{E}_{X_m}$ of derived edges of $X_m$ is defined as follows: We draw an edge between $(v , g\!\!\pmod{\Gamma^{p^m}})$ and $(v', g'\!\!\pmod{\Gamma^{p^m}})$ if there is an edge $e$ from $v$ to~$v'$ in $\mathbf{E}_X$ such that ${g'\pmod{\Gamma^{p^m}}=g \cdot \alpha_m(e) \pmod{\Gamma^{p^m}}}$. 
	
	We abbreviate the vertex $(v, 1)$ of the graph $X_m$ to $v$.  
\end{defi} 
We have a natural action of $\Gamma/\Gamma^{p^m}$ on $X_m$ given by
\[(g'\!\!\!\!\pmod{\Gamma^{p^m}}) \circ (x,g\!\!\!\!\pmod{\Gamma^{p^m}})=(x,g \cdot g'\!\!\!\!\pmod{\Gamma^{p^m}}),\] 
where $x$ is a vertex of $X$. 
By using the natural projection $\Gamma\longrightarrow \Gamma/\Gamma^{p^m}$, we also have a natural action of $\Gamma$ on $X_m$. This induces a well-defined action of $\Gamma$ on $J(X_m)$. We also obtain an action of the two canonical group rings 
\[ R = \Z_p[\Gamma], \quad \Z_p[[\Gamma]]. \] 
The completed group ring is isomorphic non-canonically to the Iwasawa algebra ${\Lambda = \Lambda_l = \Z_p[[T_1, \ldots, T_l]]}$, as has been explained in the previous subsection. 
\begin{defi} 
	We let $X_\infty$ be a $\Z_p^l$-voltage cover of $X$. Then we define 
	\[ J_\Lambda := J(X_\infty) \otimes_R \Lambda. \] 
\end{defi} 
Note that $J_\Lambda$ is a finitely generated $\Lambda$-module, i.e. it is an Iwasawa module in the language of Section~\ref{section:notation_iw-modules} (this will follow from Lemma~\ref{lemma:projective-limit} and the results in Section~\ref{section:J(X_m)}). Therefore it makes sense to speak of the (generalised) Iwasawa invariants of the graph~$X_\infty$. 

In this article, we always assume that all the graphs $X_m$ are connected. We will provide a sufficient criterion for this property below (see Lemma~\ref{lemma:connected}). Thus, $X_m/X$ is a Galois cover with group ${\Gamma/\Gamma^{p^m} \cong (\Z/p^m\Z)^l}$. Then as a $\Z_p[\Gal(X_m/X)]$-module $\textup{Pr}(X_m)$ is generated by the elements
\begin{eqnarray*} p_{v,0} & = & -\sum_{e \in \textbf{E}_v(X_m)}(t(e)-v) \\ 
	& = & -\sum_{e\in \textbf{E}_v(X),t(e)\neq v}(\alpha_m(e)t(e)-v)-\sum_{e\in \textbf{E}_v(X),t(e)=v}(\alpha_m(e)+\alpha_m(e)^{-1}-2)v,\end{eqnarray*} 
where ${v \in V(X)}$, $\textbf{E}_v(Y)$ denotes the set of directed edges with origin $v$ in a graph $Y$, and where $t(e)$ is the target of an edge $e$.

If $e$ is a loop with trivial voltage assignment then the term ${\alpha_m(e) + \alpha_m(e)^{-1} - 2}$ vanishes in the above sum. The coefficient of ${v = (v,1)}$ in $p_{v,0}$ (over $\Z_p$) is given by \[\deg(v)v-2\vert \{e\mid t(e)=v,\alpha_m(e)=1\}\vert.\] 

Now suppose that $X$ is a finite graph, and let ${n = |V(X)|}$. 
Fixing an indexing of the vertices of $X$, say ${V(X) = \{ v_1, \ldots, v_n\}}$, we will also write ${p_{i,0} = p_{v_i,0}}$ for these elements. We define a $\Z_p$-linear map
\[\mathcal{L}_\alpha\colon \textup{Div}(X_\infty)\longrightarrow  \textup{Pr}(X_\infty), \quad (v_{i},1) \mapsto p_{i,0}.\]
It is easy to see from the definitions and the action of $\Gal(X_\infty/X)$ on $\textup{Div}(X_\infty)$ that this map is unique and well-defined. It is called the \emph{voltage $p$-Laplacian}, see also~\cite{diss-gonet}. 

\begin{defi}
	We say that a graph $X$ contains a non-trivial cycle of length $n>1$ if there exist a sequence of vertices $v_0,v_1,\dots v_{n-1},v_n,v_0$ and a sequence of directed edges $e_0,\ldots, e_{n}$ such that $e_i$ is an edge from $v_i$ to $v_{i+1}$ for $0\le i\le n-1$ and $e_n$ is an edge from $v_n$ to $v_0$. We assume furthermore that $e_{i+1}\neq {e_i}^{-1} $ for ${0\le i\le n-1}$ and $e_n\neq e_0^{-1}$.
	
	For any such cycle $C$ we write $\beta_C$ for the group element $\prod_{i=0}^n\alpha(e_i)$.
\end{defi}
\begin{lemma} \label{lemma:connected} 
	$X_m$ is connected for all $m$ if and only if we can find cycles $C_1,\dots , C_l$ such that $\{\beta_{C_1},\dots,\beta_{C_l} \}$ is a set of topological generators of $\Gamma$. 
\end{lemma}
\begin{proof}
	Recall that we write $V(X_m)$ for the set of vertices of $X_m$. We say that two vertices $(v,\gamma)$ and $(v',\gamma')$ of $X_m$ are connected if there is a path connecting $(v,\gamma)$ and $(v',\gamma')$. It is easy to check that this defines an equivalence relation on $V(X_m)$. 
	
	Assume first that there are cycles $C_1,\dots, C_l$ satisfying the above condition. In what follows, we abbreviate $\beta_{C_i}$ to $\beta_i$ for each ${1 \le i \le l}$. 
	Let $v\in V(X)$. We will first show that $(v,1)$ is connected to $(v,\gamma)$ for every ${\gamma\in \Gamma/\Gamma^{p^m} \setminus \{1\}}$. As ${\beta_1,\dots ,\beta_l}$ are topological generators of $\Gamma$, it suffices to consider ${\gamma=\beta_i^{a_i}}$ for $1\le i\le l$ and ${a_i\in \Z/p^m\Z}$. Let $v_0$ be a vertex lying on $C_i$. Then ${(v_0,\gamma')}$ is connected to ${(v_0,\gamma'\cdot \beta_i^{a_i})}$ for every choice of $\gamma'$. As $X$ is connected we can find a path ${e_0,e_1,\ldots, e_k}$ from $v$ to~$v_0$ and we see that $(v,1)$ is connected to ${(v_0,\prod_{j=1}^k\alpha_m(e_j))}$ and ${(v_0,\tilde{\gamma})}$ is connected to~${(v,\tilde{\gamma}\prod_{j=1}^k\alpha_m(e^{-1}_j))}$ for every choice of $\tilde{\gamma}$. We obtain the following relations
	\[\left(v,1\right)\sim \left(v_0,\prod_{j=1}^k \alpha_m(e_j)\right)\sim \left(v_0,\prod_{j=1}^k\alpha_m(e_j)\beta^{a_i}_i\right)\sim \left(v,\beta_i^{a_i}\right)\]
	proving our first claim. 
	
	For every $v'\in V(X)$ we can find an element $\gamma\in \Gamma/\Gamma^{p^m}$ such that $(v,1)$ is connected to $(v',\gamma)$ -- using that $X$ is connected. Using the first part of the proof we see that $(v,1)\sim (v',\gamma)\sim (v',\gamma')$ for every choice of $\gamma'$. Therefore $X_m$ is indeed connected. 
	
	Assume now that all the $X_m$ are connected. In particular $X_1$ is connected. Let $v\in V(X)$ and $\gamma\in \Gamma/\Gamma^p\setminus\{1\}$ be arbitrary. Then $(v,1)$ is connected to $(v,\gamma)$. Thus, there is a path $e_1,\dots, e_k$ from $v$ to $v$ such that $\gamma=\prod_{j=1}^k\alpha_1(e_j)$. Without loss of generality we can assume that $e_i\neq e^{-1}_{i+1}$ for $1\le i\le k-1$. Let $v_i$ be the target of $e_i$ and set $v=v_0$. Let $k_v$ be the minimal index such that $e_{k_v}\neq e^{-1}_{k-k_v+1}$. Such an index exists, as ${\gamma\neq 1}$. Then $e_{k_v},\dots ,e_{k-k_v+1}$ is a path from $v_{k_v-1}$ to itself and it is indeed a cycle. Furthermore, ${\gamma=\prod_{j=k_v}^{k-k_v+1}\alpha_1(e_j)}$. 
\end{proof} 

\emph{In the rest of the paper we will always assume that all the graphs $X_m$ are connected.}


\section{The module theoretic perspective in the multidimensional case} \label{section:J(X_m)} 
In this section we generalise the module-theoretic point of view from~\cite{diss-gonet} to the case of arbitrary $l$: we will explain the structure of the Iwasawa module $\varprojlim_m J(X_m)$ and its relation to the asymptotic growth of $J(X_m)$. 


Let ${R = \Z_p[\Gamma]}$, let $\Lambda=\Z_p[[T_1,\ldots ,T_l]]$ be as before, and let $I_m$ and ${R_m = \Lambda/I_m}$ be defined as in Section~\ref{section:notation_iw-modules}. We write $\textup{Div}_\Lambda$, $\textup{Div}^0_\Lambda$, $\textup{Pr}_\Lambda$ and $\textup{Pic}_\Lambda$ for the corresponding $\Lambda$-modules $\textup{Div}(X_\infty)\otimes \Lambda$, $\textup{Div}^0(X_\infty)\otimes \Lambda$, $\textup{Pr}(X_\infty)\otimes \Lambda$ and $\textup{Pic}(X_\infty) \otimes \Lambda$. In what follows we write ${|V(X)| = n}$. It is easy to see that $\textup{Div}_\Lambda$ is $\Lambda$-free of rank $n$. 
\begin{rema}
	The ring $\Lambda$ is actually flat over $R$. The argument boils down to showing that $I\otimes \Lambda=I\Lambda$ for every ideal $I\subseteq R$. This argument involves several diagram chases and the fact that projective limits and tensor products commute under certain conditions. As the proof is not very enlightening and the computation of the necessary $\textup{Tor}$ modules is rather easy in the cases of interest to us, we decided not to include here a proof of the fact that $\Lambda$ is flat over $R$.
\end{rema}
\begin{lemma}
	\label{pr-free}
	The $\Lambda$-module $\textup{Pr}_\Lambda$ is free of rank $n$. Furthermore, the natural map $\textup{Pr}_\Lambda\longrightarrow \textup{Div}_\Lambda$ is injective. Thus, we obtain a short exact sequence
	\[0\longrightarrow \textup{Pr}_\Lambda\longrightarrow \textup{Div}_\Lambda\longrightarrow \textup{Pic}_\Lambda\longrightarrow 0.\]
\end{lemma}
\begin{proof}
	Let $V$ be the submodule of $\textup{Div}_\Lambda$ generated by the $p_{i,0}$, i.e. $V$ is the image of $\textup{Pr}_\Lambda$ in $\textup{Div}_\Lambda$. As $\textup{Div}(X_m)$ is generated by the $v_i$ as an $R_m$-module, there is clearly a surjective map 
	\[\textup{Div}_\Lambda/I_m\longrightarrow \textup{Div}(X_m).\]
	Assume that there are elements $F_i\in \Lambda$, not all zero, such that ${\sum F_i \cdot p_{i,0}=0}$ in $\textup{Div}_\Lambda$. Without loss of generality we can assume that $F_1\neq 0$. Then there is a surjection ${\Lambda/F_1\oplus \Lambda ^{n-1}\twoheadrightarrow V}$. Taking quotients by $(I_m\Lambda)^n$ we see that the $\Z_p$-rank of $\textup{Pr}(X_m)$ is at most ${(n-1)p^{ml}+E(m)}$ for some error term ${E(m) = O(p^{m(l-1)})}$. It follows that \[np^{ml}-1=\Z_p\textup{-rank}(\textup{Pr}(X_m)) \le (n-1)p^{ml}+ E(m), \quad E(m) = O(p^{m(l-1)}),\] yielding a contradiction. Thus, $V$ is $\Lambda$-free as well and ${\textup{Pr}_\Lambda\longrightarrow \textup{Div}_\Lambda}$ is injective.
	
	The exactness of the sequence in the statement of the lemma now follows from tensoring 
	\[0\longrightarrow \textup{Pr}(X_\infty)\longrightarrow\textup{Div}(X_\infty)\longrightarrow \textup{Pic}(X_\infty)\longrightarrow 0\] 
	with $\Lambda$. Note that the sequence stays left exact after tensoring with $\Lambda$ as $\textup{Pr}_\Lambda$ embeds into $\textup{Div}_\Lambda$.\end{proof}
Let $M_\Lambda \supseteq \textup{Pr}_\Lambda$ be the submodule of $\textup{Div}_\Lambda$ generated by the elements $p_{i,0}$, the elements ${v_i-v_{j}}$ for ${1 \le j < i \le n}$ and $(T_1,\ldots, T_l)\textup{Div}_\Lambda$. The following theorem generalises results of ~\cite[Section~5]{diss-gonet}. 
\begin{thm} \label{thm:jacobian} 
	We have $J(X_m)\cong M_\Lambda/(I_m \cdot \textup{Div}_\Lambda+\textup{Pr}_\Lambda)$.
\end{thm}
\begin{proof}
	Let $\pi_m\colon \textup{Div}(X_\infty)\longrightarrow \textup{Div}(X_m)$ be the canonical map of voltage covers. Note that we can see this as an $R$-module homomorphism. We start by the following series of observations:
	\begin{compactenum}[(a)]
		\item $\pi_m$ is surjective and induces surjective maps $\textup{Div}^0(X_\infty)\longrightarrow \textup{Div}^0(X_m)$ and ${\textup{Pr}(X_\infty)\longrightarrow \textup{Pr}(X_m)}$.
		\item The kernel of $\pi_m$ is given by $I_m\textup{Div}_\Lambda\cap \textup{Div}(X_\infty)$.
		\item We have $I_m\textup{Div}_\Lambda+\textup{Div}^0(X_\infty)=M_\Lambda$.
		\item $I_m\textup{Div}_\Lambda+\textup{Pr}(X_\infty)=\textup{Pr}_\Lambda+I_m\textup{Div}_\Lambda$. 
	\end{compactenum}
	Let us first see that this series of observations suffices to prove the theorem.
	By points (a) and (b) we have an isomorphism
	\[\textup{Div}^0(X_m)\cong \textup{Div}^0(X_\infty)/(I_m \textup{Div}_\Lambda\cap \textup{Div}^0(X_\infty)),\] where the right hand term is isomorphic to \[(\textup{Div}^0(X_\infty)+I_m\textup{Div}_\Lambda)/I_m\textup{Div}_\Lambda.\]
	Combining this with (c) and (d), we have
	\begin{align*}
		J(X_m)& = \textup{Div}^0(X_m) / \textup{Pr}(X_m) \\ 
		&=M_\Lambda/(\textup{Pr}_\Lambda+I_m\textup{Div}_\Lambda).
	\end{align*}
	It remains to check the observations (a)-(d). For observation (a) note that $v_i$, ${1\le i\le n}$, is a set of generators for $X_m$ (seen as an $R$-module). 
	To see that $\pi_m$ remains surjective when restricted to $\textup{Div}^0(X_\infty)$, it suffices to note that the elements ${v_{j}-v_{i}}$ and ${(v_j, \tau_k)-(v_{j},1)}$ for a  set of topological generators ${\tau_1,\ldots, \tau_l}$ of $\Gal(X_\infty/X)$ generate $\textup{Div}^0(X_m)$ as well as $\textup{Div}^0(X_\infty)$ as $R$-modules. Using a similar argument for the generators $p_{i,0}$ instead shows the corresponding claim for $\textup{Pr}(X_\infty)$ and $\textup{Pr}(X_m)$.
	
	For point (b) note that $\textup{Div}_\Lambda$ is a free $\Lambda$-module of rank $n$ with generators $(v_j,1)$ for ${1\le j\le n}$. Furthermore, $\textup{Div}(X_m)$ is a free $R_m$-module with the same generators. As $\textup{Div}(X_\infty)$ contains the $R$-module generated by $(v_j,1)$, we see that $\textup{Div}(X_\infty)$ lies densely in $\textup{Div}_\Lambda$. We obtain the following isomorphisms. 
	\begin{align*}
		\textup{Div}(X_m)&\cong \bigoplus_{i=1}^n R_m v_i\cong 
		\bigoplus_{i=1}^n \Lambda/I_m v_i\\ 
		&\cong \textup{Div}_\Lambda/I_m\\
		&\cong \textup{Div}(X_\infty)/(\textup{Div}(X_\infty)\cap I_m\textup{Div}_\Lambda),\end{align*} which proves claim (b).
	
	Now we prove (c). By construction $\textup{Div}^0(X_\infty)\subseteq M_\Lambda$, and as our chosen generators of $M_\Lambda$ lie in $\textup{Div}^0(X_\infty)$, we see that $\textup{Div}^0(X_\infty)$ lies densely in $M_\Lambda$. Note that ${I_m\textup{Div}_\Lambda}$ 
	is a neighbourhood of $0$ in $M_\Lambda$. Let $y\in M_\Lambda$ be an arbitrary element, then there exists $z\in \textup{Div}^0(X_\infty)\cap (y+I_m\textup{Div}_\Lambda)$. 
	Thus, 
	\begin{align*} y+I_m\textup{Div}_\Lambda & \subseteq z+I_m\textup{Div}_\Lambda \\ 
		& \subseteq \textup{Div}^0(X_\infty)+I_m\textup{Div}_\Lambda.\end{align*} 
	Varying $y$ in $M_\Lambda$ gives the claim.
	
	Claim (d) can be proved similarly to point (c) by using the generators $p_{i,0}$. Clearly, ${I_m\textup{Div}_\Lambda+\textup{Pr}(X_\infty)\subseteq \textup{Pr}_\Lambda+I_m\textup{Div}_\Lambda}$. By definition, $\textup{Pr}(X_\infty)$ lies densely in $\textup{Pr}_\Lambda$ and $I_m\textup{Div}_\Lambda\cap \textup{Pr}_\Lambda$ is a neighbourhood of $0\in \textup{Pr}_\Lambda$ in the subspace topology. Let now $y\in \textup{Pr}_\Lambda$. Then there exists $z\in \textup{Pr}(X_\infty) \cap (y+(I_m\textup{Div}_\Lambda\cap \textup{Pr}_\Lambda))$. Thus, 
	\[y+I_m\textup{Div}_\Lambda\subseteq \textup{Pr}(X_\infty)+I_m\textup{Div}_\Lambda.\]
	Varying $y\in \textup{Pr}_\Lambda$ gives the claim. 
\end{proof}
\begin{lemma} \label{lemma:NundL} 
	Let $L$ be a finitely generated $\Lambda$-torsion module and let $N\subseteq L$ be such that $L_m:=I_mL\subseteq N$ for each ${m \in \N}$. Assume that $N/L_m$ is finite for all $m$. Then there are non-negative integers $m_0$ and $l_0$ such that
	\begin{align} v_p(\vert N/L_m\vert) = m_0 p^{ml}+l_0 m p^{m(l-1)}+O(p^{m(l-1)}). \label{eq:star} \end{align} 
\end{lemma}
\begin{proof}
	Since $N/L_m$ is finite for each $m$, the quotient $L_0/L_m$ is finite for all $m$. In particular, $L/L_m$ and $L/L_{0}$ have the same $\Z_p$-rank. 
	
	To compute $v_p(\vert N/L_m \vert)$ note that we have
	\[\vert N/L_m\vert = \vert N/L_0\vert \cdot \vert L_0/L_m\vert. \]
	So it remains to compute $\vert L_0/L_m\vert$. 
	Since $L/L_m$ and $L/L_0$ have the same $\Z_p$-rank and therefore the kernel of the natural surjective map ${L/L_m \longrightarrow L/L_0}$ is contained in the $p$-torsion of $L/L_m$, we deduce
	\[\vert L_0/L_m\vert = \frac{\vert (L/L_m)[p^\infty]\vert }{\vert (L/L_0)[p^\infty] \vert}.\]  
	The desired result follows directly from ~\cite[Theorem~3.4]{cuoco-monsky} (note that in our case we even have $\Z_p\textup{-rank}(L/L_m)=O(1)$ instead of $O(p^{m(l-2)})$).
\end{proof}

This lemma can be applied, according to Theorem~\ref{thm:jacobian}, to the following $\Lambda$-modules: consider ${N = M_\Lambda/\textup{Pr}_\Lambda}$ and ${L = \textup{Div}_\Lambda/\textup{Pr}_\Lambda}$. By definition $v_i-v_j\in M_\Lambda$ for all $1\le i\le j\le n$. Hence, $\textup{Div}_\Lambda/M_\Lambda$ is cyclic as $\Lambda$-module. For example, we can write ${\textup{Div}_\Lambda/M_\Lambda = \Lambda v_1}$. Furthermore, $\textup{Div}_\Lambda/M_\Lambda$ is annihilated by $(T_1,\dots, T_l)$. Thus,  ${L/N\cong \textup{Div}_\Lambda/M_\Lambda = \Lambda v_1 = \Z_p v_1 \cong \Z_p}$. Furthermore, 
\[ L_m = I_m \cdot L \] 
is contained in $N$ for each ${m \in \N}$, and 
\[ N/L_m \cong J(X_m) \] 
is finite for each $m$ by Theorem~\ref{thm:jacobian}. As $L/N\cong \Z_p$ it follows that $L/L_m$ has $\Z_p$-rank~$1$ for each $m$. In particular, $L$ has to be $\Lambda$-torsion. As $N$ is a submodule of $L$ it has to be torsion as well.

In this case, we obtain the following arithmetic description of the numbers $m_0$ and~$l_0$ from~\eqref{eq:star}: 
\begin{coro} \label{cor:zusammenhangzuN} 
	If $l\ge 2$, then the parameters $m_0$ and $l_0$ from~\eqref{eq:star} are the generalised Iwasawa invariants of $L$. Since ${L/N \cong \Z_p}$ is pseudo-null over $\Lambda$, these are also the generalised Iwasawa invariants of $N$. 
	
	If $l=1$, then we obtain the Iwasawa invariants of $N$. In particular, $m_0$ equals the $\mu$-invariant of $L$, while ${l_0=\lambda(N) = \lambda(L)-1}$ in this case.
\end{coro}
\begin{proof}
	If $l\ge 2$ this follows directly from the above proof. In the case $l=1$ it is a straight forward computation over one-dimensional Iwasawa algebras (note that ${\varprojlim_m N/I_m \cong N}$, since ${\bigcap_m I_m = \{0\}}$).
\end{proof} 

\section{Computation of the order of \texorpdfstring{$J(X_m)$}{J(Xm)}} \label{section:linearalgebra} 
Let $X$ be a finite graph, and let $\alpha \colon S\longrightarrow \Gamma$ be as in Section~\ref{section:notation}.
Recall the definition of the $\Z_p$-linear map $\mathcal{L}_\alpha$ from Section~\ref{section:notation}. We can extend naturally the linear map $\mathcal{L}_\alpha$ to a $\Lambda$-linear map.
Recall that $\textup{Div}_\Lambda\cong \Lambda^n$ with basis $\{ v_1, \ldots, v_n\}$. We would like to give a matrix-representation of $\mathcal{L}_\alpha$ with respect to this basis. Let $D$ be the degree matrix of the base graph $X$ and define the voltage assignment matrix $A_\alpha=(\alpha_{i,j})_{1\le i,j\le n}$ as follows: 
\[\alpha_{i,j}=\begin{cases}
	\sum _{ e \in E(v_i,v_j)}\alpha(e) \quad & \textup{if } i\neq j,\\
	\sum_{e \in E(v_i,v_i)}\alpha(e)+\alpha(e)^{-1} \quad & \textup{if } i=j.
\end{cases}\] 
Then $\alpha_{i,j} \in \Lambda$ for all $i,j$. It is easy to verify that $\mathcal{L}_\alpha$ is represented by the transpose $\Delta_\infty^t$ of the matrix ${\Delta_\infty=D-A_\alpha}$ (here we use the $v_i$ as a basis). Recall that 
\[p_{v_i,0}=-\sum_{e\in E_{v_i}(X_\infty)}(\alpha(e)t(e)-v_i)=\deg(v_i)v_i-\sum_{j=1}^n \alpha_{i,j}v_j,\]
which is represented by the $i$-th row of $D-A_\alpha$ and not by the $i$-th column. Note that this ambiguity does not occur if we understand $\mathcal{L}_\alpha$ only as a $\Z_p$-linear map. In this case a $\Z_p$-basis of $\textup{Div}(X_m)$ is $\{(v_i,g)\mid v_i\in V(X), g\in (\Z/p^m\Z)^l\}$,  $\mathcal{L}_\alpha$ is represented by a symmetric matrix and we do not have to consider the transpose.

In what follows, we fix some ${m \in \N}$. Recall the definition of ${R_m = \Lambda/I_m}$ from Section~\ref{section:notation_iw-modules}. If we define $\Delta^t_m$ to be the representing matrix for the map 
\[ \mathcal{L}_{\alpha_m} \colon \textup{Div}(X_m) \longrightarrow \textup{Div}(X_m) \]  
(again in the $R_m$-basis  $v_i$), we obtain that $\Delta_m$ is the image of $\Delta_\infty$ in $\textup{Mat}_{n,n}(R_m)$.

Let $\Omega$ be the group of all $p$-power roots of unity in some fixed algebraic closure $\overline{\Q_p}$ of $\Q_p$. We say that two elements in $\Omega^l$ are equivalent if they are Galois conjugate to each other.  Let $\zeta=(\zeta_1,\dots,\zeta_l)\in \Omega^l$. We denote by $\Z_p[\zeta]$ the ring $\Z_p[\zeta_1,\dots,\zeta_l]$ obtained by adjoining all the components $\zeta_i$ to $\Z_p$. Following~\cite{cuoco-monsky} we define a map
\[\Lambda/I_m\longrightarrow \bigoplus \Z_p[\zeta], \quad F(T_1,\dots,T_l)\mapsto (F(\zeta_1-1,\dots, \zeta_l-1))_\zeta, \]
where the sum on the right hand side runs over a set of representatives ${\zeta = (\zeta_1, \ldots, \zeta_l)}$ for the equivalence classes in $\Omega^l[p^m]$. 
Using that $\textup{Div}_\Lambda$ is $\Lambda$-free of rank $n$, we can extend the above map linearly and define 
\[\phi_m\colon \textup{Div}_\Lambda/I_m\longrightarrow (\bigoplus \Z_p[\zeta])^n. \] 
Let $\tilde{\Lambda}_m$ be the image of $\phi_m$. Note that by ~\cite[Theorem~2.2]{cuoco-monsky}, $\phi_m$ has a finite cokernel and trivial kernel. Recall that $\textup{Pr}_\Lambda=\mathcal{L}_\alpha(\textup{Div}_\Lambda)$ and that $\Delta^t_\infty$ is a matrix representing the map $\mathcal{L}_\alpha$. 

It is immediate that $\mathcal{L}_\alpha(I_m\textup{Div}_\Lambda)\subseteq I_m\textup{Div}_\Lambda$. So $\mathcal{L}_\alpha$ induces well-defined maps on $\textup{Div}_\Lambda/I_m$. By abuse of notation, we denote these maps by $\mathcal{L}_\alpha$ again, and we define a map $\tilde{\mathcal{L}_\alpha}$ on $\tilde{\Lambda}_m$ such that
\[\phi_m\circ \mathcal{L}_\alpha=\tilde{\mathcal{L}_\alpha}\circ \phi_m\] 
on $\textup{Div}_\Lambda/I_m$. 
Note that $\tilde{\mathcal{L}_\alpha}$ is a linear map of $\Z_p$-modules. 
Recall that $\phi_m$ has finite cokernel. So we can extend $\tilde{\mathcal{L}_\alpha}$ to a linear map of $\tilde{\Lambda}_m\otimes \Q_p=(\oplus_{z=1}^k\Q_p[\zeta_z])^n$, where $\zeta_1, \ldots, \zeta_k$ represent the different equivalence classes in $\Omega^l[p^m]$ modulo Galois equivalence.

For every matrix $A$ with coefficients ${a_{i,j} \in \Lambda}$ and each ${\zeta \in \Omega^l}$ we denote by $A(\zeta - 1)$ the matrix with coefficients ${a_{i,j}(\zeta-1) \in \Z_p[\zeta]}$ (i.e. we replace $T_i$ by the $i$-th component of $\zeta - 1$). 
\begin{lemma} \label{lemma:Am} 
	We can choose a basis of $\tilde{\Lambda}_m\otimes \Q_p$ such that $\tilde{\mathcal{L}_\alpha}$ is represented by a $p^{ml}n\times p^{ml}n$-block matrix
	\[A_m= \begin{pmatrix}
		\Delta^t_\infty(\zeta_1-1)&0&\dots &\dots&\dots&0\\
		0& \Delta^t_\infty(\zeta_1-1)&\dots &\dots&\dots&0\\
		\dots&\dots&\dots&\dots&\dots&\dots\\
		0&0&\dots&\Delta^t_\infty(\zeta_2-1)&\dots&0\\
		\dots&\dots&\dots&\dots&\dots&\dots\\
		0&\dots&\dots&\dots&0&\Delta^t_\infty(\zeta_{k}-1)
	\end{pmatrix},\]
	where $\zeta_1,\ldots ,\zeta_k$ represent the different equivalence classes in $\Omega^l[p^m]$ modulo Galois equivalence and each block $\Delta_\infty^t(\zeta_i-1)$ occurs $\varphi(\textup{ord}(\zeta_i))$ times, where $\varphi$ denotes Euler's totient function. 
\end{lemma}
\begin{proof}
	Let $F$ be an element in $\textup{Div}_\Lambda$. 
	Note that $\phi_m\circ \mathcal{L}_\alpha(F)$ is a vector with components $\Delta^t_\infty(\zeta_i-1)F(\zeta_i-1)$. 
	
	We fix an index $j$ ($1\le j\le k$) and let $F_jv_i\in \Lambda v_i$ be such that $\phi_m(F_jv_i)$ vanishes at all $\zeta_z-1$ for $z\neq j$ and is an element $a\in \Z_p\setminus\{0\}$ at $\zeta_j-1$. Let $G$ be an arbitrary element in $\Lambda$. Then we have
	\[\tilde{\mathcal{L}_\alpha}(\phi_m(GF_jv_i))=\phi_m(\mathcal{L}_\alpha (GF_jv_{0, i}))=\phi_m(GF_j \Delta^t_\infty v_i).\] 
	Recall that $G(\zeta-1)F_j(\zeta-1)=0$ for all $\zeta$ that do not lie in the equivalence class of~$\zeta_j$. Thus
	\[G(\zeta-1)F_j(\zeta-1)\Delta^t_\infty(\zeta-1)=0\]
	for all these $\zeta$. In particular, $\phi_m(GF_j\Delta^t_\infty v_i)$ has only one non-trivial component, namely at $\zeta_j-1$. Note that $\phi_m((\Lambda/I_m) F_jv_i)$ has finite index in $\Z_p[\zeta_j]$. 
	Using the fact that $\bigoplus_{z=1}^k\Z_p[\zeta_z]$ is $\Z_p$-free, we have shown that if we start with an element in $\tilde{\Lambda}_m$ that is non-trivial at exactly one $\zeta_j-1$, then the same holds for its image under $\tilde{\mathcal{L}_\alpha}$.
	From this we can deduce that $\tilde{\mathcal{L}_\alpha}(\Q_p[\zeta_j])^n\subseteq (\Q_p[\zeta_j])^n$.
	
	It remains to compute the corresponding matrix. Let $G_1,\ldots ,G_{\varphi(\textup{ord}(\zeta_j))-1}$ be polynomials such that $G_w(\zeta_j-1)=\zeta_j^w$ and let $G_0=1$. Note that $\phi_m(G_wF_jv_i)$ for $0\le w\le \varphi(\textup{ord}(\zeta_j))-1$ spans a sublattice of finite index in $\Z_p[\zeta_j]$. Using the above computations (with $G_w$ for $G$) we see that
	\[\tilde{\mathcal{L}_\alpha}(\zeta_j^wa\phi_m(v_{i}))=\sum_{u=1}^n\Delta^t_\infty(\zeta_j-1)\zeta_j^wa\phi_m(v_{u}).\] 
	Therefore, we see that the block $\Delta^t_\infty(\zeta_j-1)$ has to occur $\varphi(\textup{ord}(\zeta_j))$ times in the matrix representing $\tilde{\mathcal{L}_\alpha}$ on $\tilde{\Lambda}_m\otimes \Q_p$. 
\end{proof}
\begin{lemma} We have the following identity 
	\[v_p(\vert J(X_m)\vert)=\sum_{\zeta\neq 1}v_p(\det(\Delta^t_\infty(\zeta- 1)))+v_p(|J(X_0)|), \]
	where the sum runs over all elements in $\Omega^l[p^m]$ that are different from $(1,1,\dots,1)$.
\end{lemma}

\begin{proof}
	Recall that $N/I_mL$ is finite for all $m$ and that $L/I_mL$ has $\Z_p$-rank one for each $m$. By definition of $L$ we see that \[L/I_m L=\textup{Div}_\Lambda/(I_m+\textup{Pr}_\Lambda)=\textup{Div}(X_m)/\textup{Pr}(X_m)=\textup{Div}(X_m)/\Delta^t_m\textup{Div}(X_m).\] 
	Therefore the kernel of the linear map given by multiplication with $\Delta^t_m$ has $\Z_p$-rank~$1$. Since $\phi_m$ is an injection, it follows that the matrix $A_m$ introduced in Lemma~\ref{lemma:Am} has rank $p^{ml}n-1$ for all $m$. In particular, the matrix $\Delta^t_\infty(0)$ has determinant zero, but all the matrices $\Delta^t_\infty(\zeta_i-1)$ have non-zero determinant. 
	
	Fix a basis of $\tilde{\Lambda}_m$ and let $A'_m$ be the matrix describing $\tilde{\mathcal{L}_\alpha}$ in this basis. Then $A_m$ and $A'_m$ are conjugate. 
	Moreover, $\textup{Div}(X_m)$ is a free $\Z_p$-module, and 
	\[ \textup{Div}(X_m)/\Delta^t_m \textup{Div}(X_m) \cong \tilde{\Lambda}_m/A'_m \tilde{\Lambda}_m. \] 
	We are interested in $(\textup{Div}(X_m)/\Delta^t_m\textup{Div}(X_m))[p^\infty]$. Let $B_m$ and $B'_m$ be the matrices representing 
	\[\mathcal{L}_\alpha \colon (\textup{Div}_\Lambda/I_m \otimes \Q_p)/\ker(A_m)\longrightarrow (\textup{Div}_\Lambda/I_m \otimes \Q_p)/\ker(A_m)\]
	and 
	\[\tilde{\mathcal{L}_\alpha}\colon \tilde{\Lambda}_m/\ker(A'_m)\longrightarrow \tilde{\Lambda}_m/\ker(A'_m)\]
	in our chosen basis. Then $B_m$ and $B'_m$ are conjugate and we obtain 
	\[v_p(\vert(\textup{Div}(X_m)/\Delta^t_m\textup{Div}(X_m))[p^\infty]\vert)=v_p(\det(B'_m))=v_p(\det(B_m)).\]
	
	We let $\zeta_k=1$. Then $B_m$ has the form
	\[\begin{pmatrix}
		\Delta^t_\infty(\zeta_1-1)&0&\dots &\dots&\dots&0\\
		0& \Delta^t_\infty(\zeta_1-1)&\dots &\dots&\dots&0\\
		\dots&\dots&\dots&\dots&\dots&\dots\\
		0&0&\dots&\Delta^t_\infty(\zeta_2-1)&\dots&0\\
		\dots&\dots&\dots&\dots&\dots&\dots\\
		0&\dots&\dots&\dots&0&B_0
	\end{pmatrix},\]
	where $v_p(\det(B_0))= v_p(\vert J(X_0)\vert)$.
	Using the structure of the matrix $B_m$, the claim follows. 
\end{proof}

In particular, together with ~\cite[Theorem 5.6]{monsky} this proves Greenberg's conjecture, i.e. we obtain the following 
\begin{thm}
	\label{thm:greenberg-conj}
	Greenberg's conjecture holds true for the growth of $v_p(|J(X_m)|)$, i.e. there exists a polynomial ${Q(X,Y) \in \Z_p[X,Y]}$ of total degree $l$ and degree at most one in the second variable such that 
	\[ v_p(|J(X_m)|) = Q(p^m,m) \] 
	for all sufficiently large ${m \in \N}$. 
\end{thm}
\begin{rema}
	Our method gives the same polynomial as the method using Ihara $L$-functions from ~\cite[Theorem~6.2]{dubose-vallieres}. This follows directly from the construction of the polynomial $P(X,Y)$ occurring in the statement of \emph{loc.~cit.} The difference in the methods of proof lies in the fact that \emph{loc.~cit.} relates this polynomial to Ihara $L$-functions which in turn describe the number of spanning trees, whereas the proof in the present paper is based on the fact that the matrix $\Delta_\infty$ describes the Jacobian and that the cardinality of the Jacobian is precisely the number of spanning trees. 
\end{rema}




\section{An Iwasawa main conjecture} \label{section:main_conjecture} 
Let as before $X$ be a finite connected graph. Let $X_\infty$ be a $\Z_p^l$-voltage cover of $X$. For every finite graph we define the zeta function and the Ihara $L$-functions as in~\cite{dubose-vallieres}.
Let~$\Delta_\infty$ be defined as in Section~\ref{section:linearalgebra}. Let $\psi$ be a character of $\Z_p^l$ with finite image and let $A_\psi$ be the matrix defined in~\cite{dubose-vallieres}. From the definitions in~\cite{dubose-vallieres} we see that 
\[P_{\psi}(u):=\det(I-A_\psi u+(D-I)u^2)=\frac{(1-u^2)^{\chi(X)}}{L_X(u,\psi)}.\]
In this section we assume that no vertex has degree $1$. If $\psi$ is not the trivial character then $P_\psi(1)\neq 0$ (see for example ~\cite[Section 5]{vallieres3}). Using the definition of $A_\psi$, it follows that
\[P_\psi(1)=\psi(\det(\Delta_\infty)).\]
If we interpret the function $\lim_{u\to 1}\frac{(1-u^2)^{\chi(X)}}{L_X(u,\psi)}$ as the algebraic part of the Ihara $L$-function, then $\det(\Delta_\infty)\in \Lambda$ can be seen as a $p$-adic $L$-series interpolating special values of the algebraic part of $L_X(u,\psi)$. 

To understand the algebraic side we have to analyse the $\Lambda$-module structure of $\textup{Pic}_\Lambda$ and $N$ in more detail. Recall that $\textup{Pic}_\Lambda$ was defined as $\textup{Pic}(X_\infty)\otimes \Lambda$. 

\begin{lemma} \label{lemma:Mlambda} 
	\[M_\Lambda\cong \textup{Div}^0_\Lambda.\]
\end{lemma}
\begin{proof}
	Recall that $\textup{Div}^0(X_\infty)+I_m\textup{Div}_\Lambda=M_\Lambda$ for all $m$ (by observation~(c) in the proof of Theorem~\ref{thm:jacobian}). It follows that the image of $\textup{Div}^0_\Lambda$ in $\textup{Div}_\Lambda$ is equal to $M_\Lambda$. Note that $\textup{Div}(X_\infty)/\textup{Div}^0(X_\infty)$ is annihilated by $I_0=(T_1,T_2,\dots, T_l)$. 
	It follows that
	\[I_0^{-1}\textup{Tor}_1^R(\Z_p,\Lambda)=\textup{Tor}_1^R(\Z_p\otimes I_0^{-1}R,\Lambda\otimes RI_0^{-1})=0.\]
	Therefore $\textup{Tor}_1^R(\Z_p,\Lambda)$ is torsion as its annihilator intersects with $I_0$ non-trivially.
	But 
	$$\textup{Div}^0(X_m)[I_0]=(\prod_{i=1}^l\nu_{m,0}(T_i))\textup{Div}(X_m)\cap \textup{Div}^0(X_m).$$
	It follows that $\textup{Div}^0_\Lambda$ does not contain non-trivial $I_0$-torsion. In particular, the image of 
	$\textup{Tor}_1^R(\textup{Div}(X_\infty)/\textup{Div}^0(X_\infty),\Lambda)$ in $\textup{Div}^0_\Lambda$ is trivial and we obtain a natural injection 
	\[\textup{Div}^0_\Lambda\longrightarrow \textup{Div}_\Lambda.\]
	Thus, $\textup{Div}^0_\Lambda\cong M_\Lambda$. 
\end{proof}
\begin{lemma}
	\label{lemma:projective-limit}
	Let $X_m$ be the level $m$ subcover of $X_\infty/X$. Then we have
	\[J_\Lambda:=J(X_\infty)\otimes \Lambda=\varprojlim_{m}J(X_m) \cong \varprojlim_m N/(I_m L) \cong N.\] 
	In particular, $J_\Lambda$ is a finitely generated $\Lambda$-module and we have a short exact sequence
	\[0\longrightarrow \textup{Pr}_\Lambda\longrightarrow \textup{Div}^0_\Lambda\longrightarrow J_\Lambda\longrightarrow 0.\]
\end{lemma} 
\begin{proof}
	By definition $J(X_\infty)=\textup{Div}^0(X_\infty)/\textup{Pr}(X_\infty)$ as $R$-modules. Consider the exact sequence
	\begin{equation}\label{trivial-sequence}0\longrightarrow \textup{Pr}(X_\infty)\longrightarrow \textup{Div}^0(X_\infty)\longrightarrow J(X_\infty)\longrightarrow 0.\end{equation}
	
	
	Upon tensoring \eqref{trivial-sequence} with $\Lambda$ we obtain, by Lemma~\ref{lemma:Mlambda}, an exact sequence
	\[\textup{Pr}_\Lambda\longrightarrow M_\Lambda\longrightarrow J_\Lambda \longrightarrow 0.\]
	Recall that $\textup{Pr}_\Lambda=\mathcal{L}_\alpha(\textup{Div}_\Lambda)\subseteq M_\Lambda$ injects into $\textup{Div}_ \Lambda$ by Lemma~\ref{pr-free}. In particular, the natural map $\textup{Pr}_\Lambda\longrightarrow M_\Lambda$ is injective and we obtain an isomorphism
	\begin{equation}
		\label{trivial-isomorphism}
		M_\Lambda/\textup{Pr}_\Lambda \cong J_\Lambda.
	\end{equation}
	
	Taking projective limits of the exact sequence
	\[0\longrightarrow \textup{Pr}(X_m)\longrightarrow \textup{Div}^0(X_m)\longrightarrow J(X_m)\longrightarrow 0\] gives us a sequence
	\[0\longrightarrow \varprojlim_m\textup{Pr}(X_m)\longrightarrow \varprojlim_m\textup{Div}^0(X_m)\longrightarrow \varprojlim_mJ(X_m).\]
	By Theorem~\ref{thm:jacobian} $\varprojlim_mJ(X_m)\cong N$. Note that $\textup{Pr}(X_m)$ is generated by the $p_{i,0}$ and that these generators do not have a relation over $\Lambda$. It follows that $\varprojlim_m\textup{Pr}(X_m)\cong \textup{Pr}_\Lambda$. The module $\textup{Div}^0(X_m)$ is generated by $I_0\textup{Div}^0(X_m)$, $v_i-v_{j}$ for $1\le i<j\le n$ and the elements $p_{i,0}$. It follows that $\varprojlim_m\textup{Div}^0(X_m)$ is generated by the same elements over $\Lambda$ and we obtain that $\varprojlim_m\textup{Div}^0(X_m)\cong M_\Lambda$. Summarising, we obtain an exact sequence
	\[0\longrightarrow \textup{Pr}_\Lambda\longrightarrow M_\Lambda\longrightarrow \varprojlim_mJ(X_m)\longrightarrow 0. \]
	Together with \eqref{trivial-isomorphism} the claim follows.
\end{proof}


Fixing this notation, we can prove an \emph{Iwasawa main conjecture} for graphs: 
\begin{thm} \label{thm:main_conjecture} 
	Let $l\ge 2$. Then 
	\[\textup{Char}(J_\Lambda)=(\det(\Delta_\infty)). \] 
\end{thm}
\begin{proof}
	Recall that $\Delta^t_\infty$ is the matrix representing the $\Lambda$-linear map 
	\[\mathcal{L}_\alpha \colon \textup{Div}_\Lambda \longrightarrow \textup{Div}_\Lambda, \]
	such that $\mathcal{L}_\alpha(\textup{Div}_\Lambda)=\textup{Pr}_\Lambda$. 
	By Lemma~\ref{pr-free} we have a short exact sequence 
	\[0\longrightarrow \textup{Pr}_\Lambda\longrightarrow \textup{Div}_\Lambda\longrightarrow (\textup{Div}(X_\infty)/\textup{Pr}(X_\infty))\otimes \Lambda\longrightarrow 0.\]
	
	By  Lemma~\ref{lemma:projective-limit} we have a second exact sequence  
	\[0\longrightarrow \textup{Pr}_\Lambda\longrightarrow \textup{Div}^0_\Lambda\longrightarrow J_\Lambda\longrightarrow 0.\]
	
	There is a natural injection
	\[\psi\colon J_\Lambda\longrightarrow (\textup{Div}(X_\infty)/\textup{Pr}(X_\infty))\otimes \Lambda,\]
	whose cokernel has $\Z_p$-rank one. In particular
	\[\textup{Char}((\textup{Div}(X_\infty)/\textup{Pr}(X_\infty))\otimes \Lambda)=\textup{Char}(J_\Lambda), \] 
	since ${l \ge 2}$ and therefore each finitely generated $\Z_p$-module is pseudo-null over $\Lambda$. 
\end{proof}
\begin{rema} \label{rem:main_conjecture} 
	Note that this proof crucially depends on the fact that $l\ge 2$. In the case $l=1$, we see that ${\textup{Char}(\coker \psi)=(T)}$. Thus, in this case
	\[(T) \cdot \textup{Char}(J_\Lambda)= (\det(\Delta_\infty)).\]
\end{rema}
